\scanpage{40}

ou tout au moins que le premier mb. soit de torsion. En fait, il est de \uline{rang infini} sur $W$ !
\[
\begin{aligned}
\mathbb H^1(\mathfrak X,\Omega^\bullet_{\mathfrak X/S})
&=\Omega^1(\mathfrak X/S)/d(\mathcal O(\mathfrak X/S))\\
&\simeq W\{t\}/\underbrace{\operatorname{Im}(f\rightsquigarrow f')}_{\substack{\text{séries formelles}\\\text{restreintes }\uline{\text{dérivées}}}}
\end{aligned}
\]
Or pour que $\sum a_n t^n$ soit une dérivée en $\sum b_n t^n$, il f. et s. que
\[
\sum\frac1{n+1}a_n t^{n+1}
\]
\inserted{ou aussi $\sum\frac1n a_{n-1}t^n$}
soit encore une série formelle restreinte, i.e. que
\[
\unsure{v_p|\nu_p}(a_n)-\unsure{v_p|\nu_p}(n+1)\longrightarrow+\infty
\]
et comme $\unsure{v_p|\nu_p}(n)$ (quand $n$ variable) peut prendre des valeurs arbitrairement grandes, on voit qu'il ne suffit pas que
\[
\unsure{v_p|\nu_p}(a_n)\longrightarrow+\infty\quad !
\]

En notations multiplicatives, si $|\ |_{p}$ est la valeur absolue $p$-adique, \deleted{\unsure{on}} on veut que
\[
\frac{|a_n|_p}{|n+1|_p}\longrightarrow 0
\]
donc que
\[
|a_n|_p=\varepsilon_n|n+1|_p\qquad \varepsilon_n\longrightarrow 0
\]

\margin{39}
